
\section{FGLS Monad on [Ctx,Set]}
\label{sec:fgls-ctx}

There is an inclusion functor $\vJ : \worldC \inclto \ctxC$ which by precomposition extends to a functor $\vJ^* : \funC\ctxC\setC \to \funC\worldC\setC$ and further an adjunction
\begin{tikzcd}
    {\funC\ctxC\setC} \arrow[rr, "\vJ^*", shift left=3] & \perp & {\funC\worldC\setC} \arrow[ll, "\vJ_*", shift left=3]
\end{tikzcd}
by taking the right Kan extension along $\vJ$.
The idea behind $\vJ_*$ is that while locations $\vl_1, \vl_2$ in $\worldC$ represent distinct locations, variables $\va, \vb$ in $\ctxC$ represent non-distinct aliases to locations.
To recover this aliasing, $\vJ_*$ considers all the possible ways this aliasing can happen.
                                
The $\vJ_* \vdash \vJ_*$ adjunction gives us a way to extend the full-ground monad $\monFLS$ to $\ctxC$ presheaves as:
    \[ \vX \mapsto \vJ_* (\monFLS(\vX\vJ)) : \funC\ctxC\setC \to \funC\ctxC\setC \]

\noindent
We shall now give an explicit definition of this monad using the following theorem:
\begin{theorem}\label{thm:ran-coslice-ends}
    Let $\vF : \mcatC \to \mcatD, \vA,\vB : \mcatC \to \setC$.
    We will show that:
        \[ \Ran\vF{\Exp\vB\vA}(\vD) \iso \End{\vf : \vD \to \vF\vC \in \commaC\vD\vF} \Exp{\vB\vC}{\vA\vC} \]
\end{theorem}
\begin{proof}
\begin{align*}
    \Ran\vF{\Exp\vB\vA}(\vD) \iso & \End{\vC \in \mcatC} \homF\mcatD\vD{\vF\vC} \to \Exp\vB\vA\vC        & \tag{\text{\cite{kelly2022basic} 4.24}} \\
                             \iso & \End{\vf : \vD \to \vF\vC \in \commaC\vD\vF} \Exp\vB\vA\vC         & \tag{\cref{adx:coslice-ends-proof}} \\
                             \iso & \End{\vf : \vD \to \vF\vC \in \commaC\vD\vF} \Exp{\vB\vC}{\vA\vC}  & \tag{\cref{adx:coslice-end-exponent-proof}}
\end{align*}
\end{proof}

\paragraph{The Monad}

Notice that:
\[ \vJ_* (\monFLS(\vX\vJ)) = \vJ_*\vU_*(\Exp{\vP(\vX\vJ\vU)}\heapF) = \Ran{\vJ\vU}{(\Exp{\vP(\vX\vJ\vU)}\heapF)} \]
We express the FGLS monad on $\funC\ctxC\setC$ using this observation and \cref{thm:ran-coslice-ends} as:

\[\begin{array}{rcl}
    \monFLSCtx & : & \funC\ctxC\setC \to \funC\ctxC\setC \\
    \monFLSCtx\vX & \eqdef & \End{\rho : \vw \to \vw' \in \commaC\vw{\vJ\vU}} \heapF\vw' \to \Coend{\rho' : \vw' \to \vw'' \in \commaC{\vw'}\vU} (\vX\vJ\vU)(\vw'') \times \heapF\vw''
\end{array}\]

\begin{theorem}
    The $\monFLSCtx$ monad is non-finitary.
\end{theorem}
\begin{proof}
    Follows directly from the non-finitary proof of $\monFLS$.
\end{proof}

\paragraph{Naturality-style}
\label{sec:fgls-ctx-natural}

We can give an alternative description of $\monFLSCtx$ using natural transformations by observing the action of the End.
Intuitively, its action is that every $\vf_{\rho_2}$ computation can be factored through a computation $\vf_{\rho_1}$ that requires a smaller starting heap followed by an extension of the result in a compatible way.
This forbids $\vf$ from using locations that are not \emph{reachable} from $\vw$ while at the same time allowing it to be called in all compatible future worlds.

Recall that the End quantifies over all morphisms $\rho : \vw \to \vw'$, and thus gives an $\rho$-ary product: 
\[ \phi \in \Prod{\vw \xto\rho \vw' \in \objC{\commaC\vw{\vJ\vU}}} \heapF\vw' \to \vP(\vX\vJ\vU \times \heapF)\vw' \]
such that the following coherence condition holds:

for every initialisation of the form:
 \begin{tikzcd} & {\vw} \arrow[dl, "\rho_1"'] \arrow[dr, "\rho_2"] & \\ 
               {\vw_1} \arrow[rr, "\vJ\vU\varepsilon"] && {\vw_2} 
 \end{tikzcd} and heap $\vh : \heapF\vw_1$
we have: \[\vP(\vX\vJ\vU\times\heapF)(\varepsilon)(\phi_{\rho_1}(\vh)) = \phi_{\rho_2}(\heapF(\varepsilon)\vh)\]

\noindent
 If we compare this data with the data of the following natural tranformations we see that they are isomorphic:
\[\begin{array}{c}
 \End{\rho : \vw \to \vw' \in \commaC\vw{\vJ\vU}}\heapF\vw' \to \vP(\vX\vJ\vU \times \heapF)\vw' \iso
    \homF{\funC\initC\setC}{\homF\ctxC\vw{\vJ\vU(\dash)} \times \heapF}{\vP(\vX\vJ\vU \times \heapF)}
\end{array}\]

We can go further and define the functorial action, unit and join of the FGLS monad in terms of this naturality style
which will be used in \cref{sec:finitary-monad} to define the finitary monad.

\paragraph{Functorial action}

\[\begin{array}{r@{\,}c@{\,}l}
    \monFLSCtx\vX(\vf : \vw \to \vw') & : & \monFLSCtx\vX\vw \to \monFLSCtx\vX\vw' \\
    \phi & \mapsto &  \homF\ctxC{\vw'}{\vJ\vU(\dash)} \times \heapF \xto{(\dash \comp \vf) \times 1} 
                      \homF\ctxC\vw{\vJ\vU(\dash)} \times \heapF \xto\phi \vP(\vX\vJ\vU \times \heapF)
\end{array}\]

\paragraph{Unit of Monad}

\[\begin{array}{r@{\,}c@{\,}l}
    \eta_\vX & : & \vX \To \monFLSCtx\vX \\
    \vx : \vX\vw & \mapsto & \homF\ctxC\vw{\vJ\vU(\dash)} \times \heapF \xto{\vX(\dash)(\vx) \times 1} \vX\vJ\vU \times \heapF 
                                                                                          \xto{\eta^\vP_{\vX\vJ\vU \times \heapF}} \vP(\vX\vJ\vU \times \heapF)
\end{array}\]

where $\vX(\dash)(\vx) : \homF\ctxC\vw{\vJ\vU(\dash)} \To \vX\vJ\vU$ is defined as $\vf \mapsto \vX(\vf)(\vx)$.

\paragraph{Join of Monad}

\[\begin{array}{r@{\,}c@{\,}l}
    \mu_\vX & : & \Exp\monFLSCtx2\vX \to \monFLSCtx\vX \\
    \phi & \mapsto & \homF\ctxC\vw{\vJ\vU(\dash)} \times \heapF \xto\phi \vP(\monFLSCtx\vX\vJ\vU \times \heapF) \xto{\vP(\eval)} \Exp\vP2(\vX\vJ\vU \times \heapF)
    \xto{\mu^\vP_{\vX\vJ\vU \times \heapF}} \vP(\vX\vJ\vU \times \heapF)
\end{array}\]

where $\eval : (\monFLSCtx\vX)\vJ\vU \times \heapF \To \vP(\vX\vJ\vU \times \heapF)$ is defined as $\tup\phi\vh_\vw \mapsto \phi_\vw\tup{1_\vw}\vh$.

